%%%%%%%%%%%%%%%%%
% This is an sample CV template created using altacv.cls
% (v1.7.4, 30 Jul 2025) written by LianTze Lim (liantze@gmail.com), based on the
% CV created by BusinessInsider at http://www.businessinsider.my/a-sample-resume-for-marissa-mayer-2016-7/?r=US&IR=T
%
%% It may be distributed and/or modified under the
%% conditions of the LaTeX Project Public License, either version 1.3
%% of this license or (at your option) any later version.
%% The latest version of this license is in
%%    http://www.latex-project.org/lppl.txt
%% and version 1.3 or later is part of all distributions of LaTeX
%% version 2003/12/01 or later.
%%%%%%%%%%%%%%%%

%% Use the "normalphoto" option if you want a normal photo instead of cropped to a circle
% \documentclass[10pt,a4paper,withhypeper,normalphoto]{altacv}

\documentclass[10pt,a4paper,withhyper]{altacv}
%% AltaCV uses the fontawesome5 and simpleicons packages.
%% See http://texdoc.net/pkg/fontawesome5 and http://texdoc.net/pkg/simpleicons for full list of symbols.

% Change the page layout if you need to
\geometry{left=1.25cm,right=1.25cm,top=1.5cm,bottom=1.5cm,columnsep=1.2cm}

% The paracol package lets you typeset columns of text in parallel
\usepackage{paracol}


% Change the font if you want to, depending on whether
% you're using pdflatex or xelatex/lualatex
% WHEN COMPILING WITH XELATEX PLEASE USE
% xelatex -shell-escape -output-driver="xdvipdfmx -z 0" mmayer.tex
\iftutex
  % If using xelatex or lualatex:
  \setmainfont{Lato}
\else
  % If using pdflatex:
  \usepackage[default]{lato}
\fi

% Change the colours if you want to
\definecolor{VividPurple}{HTML}{3E0097}
\definecolor{SlateGrey}{HTML}{2E2E2E}
\definecolor{LightGrey}{HTML}{666666}
% \colorlet{name}{black}
% \colorlet{tagline}{PastelRed}
\colorlet{heading}{VividPurple}
\colorlet{headingrule}{VividPurple}
% \colorlet{subheading}{PastelRed}
\colorlet{accent}{VividPurple}
\colorlet{emphasis}{SlateGrey}
\colorlet{body}{LightGrey}

% Change some fonts, if necessary
% \renewcommand{\namefont}{\Huge\rmfamily\bfseries}
% \renewcommand{\personalinfofont}{\footnotesize}
% \renewcommand{\cvsectionfont}{\LARGE\rmfamily\bfseries}
% \renewcommand{\cvsubsectionfont}{\large\bfseries}

% Change the bullets for itemize and rating marker
% for \cvskill if you want to
\renewcommand{\cvItemMarker}{{\small\textbullet}}
\renewcommand{\cvRatingMarker}{\faCircle}
% ...and the markers for the date/location for \cvevent
% \renewcommand{\cvDateMarker}{\faCalendar*[regular]}
% \renewcommand{\cvLocationMarker}{\faMapMarker*}


% If your CV/résumé is in a language other than English,
% then you probably want to change these so that when you
% copy-paste from the PDF or run pdftotext, the location
% and date marker icons for \cvevent will paste as correct
% translations. For example Spanish:
% \renewcommand{\locationname}{Ubicación}
% \renewcommand{\datename}{Fecha}


%% Use (and optionally edit if necessary) this .tex if you
%% want to use an author-year reference style like APA(6)
%% for your publication list
% \input{pubs-authoryea}

%% Use (and optionally edit if necessary) this .tex if you
%% want an originally numerical reference style like IEEE
%% for your publication list
\input{pubs-num}

%% sample.bib contains your publications
\addbibresource{sample.bib}

\begin{document}
\name{Victor Hugo Nagahama}
\tagline{}
% Cropped to square from https://en.wikipedia.org/wiki/Marissa_Mayer#/media/File:Marissa_Mayer_May_2014_(cropped).jpg, CC-BY 2.0
%% You can add multiple photos on the left or right
% \photoR{2.5cm}{mmayer-wikipedia-cc-by-2_0}
% \photoL{2cm}{Yacht_High,Suitcase_High}
\personalinfo{%
  % Not all of these are required!
  \email{victorh.nagahama@gmail.com}
  \phone{(+353) 083 090 3889}
  % \mailaddress{Address, Street, 00000 County}
  \location{Maringá, Paraná, Brasil}
  \linkedin{nagahamavh}
  % \twitter{@marissamayer}
  % \xtwitter{@marissamayer}
  \homepage{nagahamavh.github.io}
  \github{nagahamavh} % I'm just making this up though.
  \orcid{0009-0008-2774-5308} % Obviously making this up too.
  %% You can add your own arbitrary detail with
  %% \printinfo{symbol}{detail}[optional hyperlink prefix]
  %% \printinfo{\faPaw}{Hey ho!}
  %% Or you can declare your own field with
  %% \NewInfoFiled{fieldname}{symbol}[optional hyperlink prefix] and use it:
  % \NewInfoField{gitlab}{\faGitlab}[https://gitlab.com/]
  % \gitlab{your_id}
	%%
  %% For services and platforms like Mastodon where there isn't a
  %% straightforward relation between the user ID/nickname and the hyperlink,
  %% you can use \printinfo directly e.g.
  % \printinfo{\faMastodon}{@username@instace}[https://instance.url/@username]
  %% But if you absolutely want to create new dedicated info fields for
  %% such platforms, then use \NewInfoField* with a star:
  % \NewInfoField*{mastodon}{\faMastodon}
  %% then you can use \mastodon, with TWO arguments where the 2nd argument is
  %% the full hyperlink.
  % \mastodon{@username@instance}{https://instance.url/@username}
}

\makecvheader

%% Depending on your tastes, you may want to make fonts of itemize environments slightly smaller
\AtBeginEnvironment{itemize}{\small}

%% Set the left/right column width ratio to 6:4.
\columnratio{0.6}

% Start a 2-column paracol. Both the left and right columns will automatically
% break across pages if things get too long.
\begin{paracol}{2}

% ------------------- Experience
\cvsection{Experiência profissional}

\cvevent{Associado (Meio período)}{Grant Thornton Irlanda}{Jan 2023 -- Ago 2023}{Dublin, Co. Dublin, Irlanda}
Planejamento e implementação de uma ferramenta interativa para problemas de previsão. Trabalhei com \textbf{modelagem de séries temporais hierarquicas}, desenvolvimento de \textbf{bibliotecas em R} e \textbf{Shiny dashboard}.

\divider

\cvevent{Estagiário em Quantitative Risk}{Grant Thornton Irlanda}{Jul 2022 -- Set 2022}{Dublin, Co. Dublin, Irlanda}
Desenvolvimento de modelos estatísticos e de machine learning para área de ESG (Environmental, Social and Governance), bem como para risco ambiental e climático. Trabalhei com \textbf{algoritmos baseado em arvores}, \textbf{modelos lineares misto} and \textbf{fluxo de ciência de dados}.

\divider

\cvevent{Residente em Inteligência Artificial}{Hub de Inteligência Artificial SENAI}{Jul 2020 -- Jul 2021}{Londrina, Paraná, Brasil}
Planjamento e implementação de soluções baseadas em modelos de machine learning e inteligência artificial para resolver problemas reais da indústria. Trabalhei com \textbf{deep learning}, \textbf{visão computational}, \textbf{computação paralela}, \textbf{web crawling} and \textbf{data mining}.

\divider

\cvevent{Cientista de dados}{Trade Technology}{Dez 2018 -- Abr 2020}{Maringá, Paraná, Brasil}
Planejamento e implementação de soluções baseadas em modelos de machine learning e inteligência artificial, além de desenvolver um sistema automatizado de relatórios em R integrado a bancos de dados SQL conjuntamente com um serviço de envio de e-mails. Atuei em todo o fluxo de ciência de dados, desde a coleta e limpeza dos dados até a análise e comunicação dos resultados. Trabalhei com \textbf{modelagem de séries temporais}, \textbf{programação linear}, \textbf{fluxo de ciência de dados}, \textbf{data mining}, \textbf{web scraping} and \textbf{visualização de dados}.

\divider

\cvevent{Estagiário}{Trade Technology}{Set 2018 --  Dez 2018}{Maringá, Paraná, Brasil}
Desenvolvimento de um sistema automatizado de relatórios em R integrado a bancos de dados SQL conjuntamente com um serviço de envio de e-mails. Atuei em todo o fluxo de ciencia de dados, desde a coleta e limpeza dos dados até a análise e comunicação dos resultados. Trabalhei com \textbf{fluxo de ciência de dados}, \textbf{data mining}, \textbf{web scraping} and \textbf{visualização de dados}.

\newpage

% ------------------- Projects
\cvsection{Projetos}

\cvevent{}{CoVideo-19: Monitoramento de Multidões em Tempo Real por meio de Vigilância por Vídeo}{Jun 2020 -- Jul 2021}{}
Desenvolvimento de um modelo de visão computacional para detectar e analisar a dinâmica de multidões em vias públicas utilizando câmeras de vídeo em tempo real. Trabalhei com \textbf{deep learning}, \textbf{visão computational} and \textbf{análise de dados espaço-temporais}.

% ------------------- Publications
\cvsection{Publicações acadêmicas}

%% Specify your last name(s) and first name(s) as given in the .bib to automatically bold your own name in the publications list.
%% One caveat: You need to write \bibnamedelima where there's a space in your name for this to work properly; or write \bibnamedelimi if you use initials in the .bib
%% You can specify multiple names, especially if you have changed your name or if you need to highlight multiple authors.
\mynames{Victor\bibnamedelima Hugo\bibnamedelima Nagahama}
  
%% MAKE SURE THERE IS NO SPACE AFTER THE FINAL NAME IN YOUR \mynames LIST

\nocite{*}

\printbibliography[heading=pubtype,title={\printinfo{\faFile*[regular]}{Artigos Acadêmicos}}, type=misc]

%\divider
%\printbibliography[heading=pubtype,title={\printinfo{\faUsers}{Presentations}}, type=inproceedings]

% ----------------------------- Column 2
\switchcolumn

\cvsection{Idiomas}

\cvskill{Português}{5}
\cvskill{Inglês}{4.5}

% ------------------- Technologies
\cvsection{Tecnologias}

\cvskill{R}{5}
\cvskill{Python}{5}
\cvskill{Git}{5}
\cvskill{Stan}{4.5}
\cvskill{SQL}{4}
\cvskill{Linux}{4}
\cvskill{SAS}{4}
\cvskill{Docker}{3.5}
\cvskill{NoSQL}{3}
\cvskill{C++}{3}

% ------------------- Education
\cvsection{Formação}

\cvevent{Doutorado em Ciência de dados}{Maynooth University}{Set 2021 -- atualmente}{}

\divider
\cvevent{Pós Graduação em Inteligência Artificial aplicado a Industria}{Faculdade de Tecnologia SENAI}{Ago 2020 -- Jul 2021}{}

\divider
\cvevent{Bacharel em Estatística}{Universidade estadual de Maringá}{Fev 2014 -- Fev 2019}{}

% ------------------- Strengths
\cvsection{Habilidades}

\cvtag{Modelagem estatística}
\cvtag{Machine Learning}
\cvtag{Inteligência Artificial}
\cvtag{Ciência de dados}

\divider \smallskip

\cvtag{Pensamento critico}
\cvtag{Inovação}
\cvtag{Autodidata}
\cvtag{Colaboração}

\newpage

% ------------------- Teaching experience
\cvsection{Experiência acadêmica}

\cvevent{Tutor}{Maynooth University}{2022 -- 2024}{}
\begin{itemize}
    \item Introduction to Data Science 1 (DS151)
    \item Introduction to Data Science 2 (DS152)
    \item Data Analysis (ST201)
    \item Design and Analysis of Experiments (ST202)
    \item R for Statistics and Data Science (ST203)
\end{itemize}

\divider
\cvevent{Tutor}{Universidade Estadual de Maringá}{2015 -- 2016}{}
\begin{itemize}
    \item Probabilidade e Inferência estatística
    \item Probabilidade e estatística
\end{itemize}

\divider
\cvevent{Instrutor de mini curso}{Universidade Estadual de Maringá}{2015}{}
\begin{itemize}
    \item Introdução ao R
\end{itemize}

% ------------------- Others
\cvsection{Outros}

{\large \faIcon{briefcase} \textbf{Empresa Junior}}

\begin{tabular}{{p{0.11\textwidth}p{0.25\textwidth}}}
	2016 - 2017 & \textbf{Presidente} \newline
	Estats Consultoria \\
	2015 - 2016 & \textbf{Diretor de comunicação} \newline
	Estats Consultoria \\
	2014 - 2015 & \textbf{Assistente de projetos} \newline
	Estats Consultoria \\
\end{tabular}

\end{paracol}

\end{document}
